\par\clearpage
\SourceLeaf{142}\GaussVTheoryEarlyFolio{138}
\[
=\frac{(1-\varepsilon)^2R\cdot\mathrm{d}u}{(1-(2\varepsilon-\varepsilon\varepsilon)\cos u^2)^{\frac32}}
\]
und die Grösse eines Elements des Parallelkreises
\[
=\frac{R\sin u\cdot\mathrm{d}\lambda}{\BareRoot{(1-(2\varepsilon-\varepsilon\varepsilon)\cos u^2)}}
\]

Zerlegt man die horizontale magnetische Kraft in zwei Theile, wovon der eine $X$ in der Richtung des Erdmeridians, der andere $Y$ senkrecht dagegen wirkt, und betrachtet man als positiv $X$, insofern diese Componente nach Norden, und $Y$, insofern diese nach Westen gerichtet ist, so wird
\[
\begin{aligned}
X&=-\frac{(1-(2\varepsilon-\varepsilon\varepsilon)\cos u^2)^{\frac32}}{(1-\varepsilon)^2}\cdot\frac{\mathrm{d}V}{R\mathrm{d}u}\\
Y&=-\BareRoot{(1-(2\varepsilon-\varepsilon\varepsilon)\cos u^2)}\cdot\frac{\mathrm{d}V}{R\sin u\cdot\mathrm{d}\lambda}
\end{aligned}
\]
Die ganze horizontale Kraft wird sodann
\[
=\BareRoot{(XX+YY)}
\]
und die Tangente der Declination
\[
=\frac YX
\]

Vernachlässigt man das Quadrat der Abplattung $\varepsilon$, so werden jene Ausdrücke
\[
\begin{aligned}
X&=-(1+(2-3\cos u^2)\varepsilon)\cdot\frac{\mathrm{d}V}{R\mathrm{d}u}\\
Y&=-(1-\varepsilon\cos u^2)\cdot\frac{\mathrm{d}V}{R\sin u\cdot\mathrm{d}\lambda}
\end{aligned}
\]
oder wenn man die Abplattung ganz bei Seite setzt
\[
X=-\frac{\mathrm{d}V}{R\mathrm{d}u}\qquad Y=-\frac{\mathrm{d}V}{R\sin u\cdot\mathrm{d}\lambda}
\]

Die bis jetzt zu Gebote stehenden Beobachtungsdata sind noch viel zu dürftig, und die meisten derselben viel zu roh, als dass es gegenwärtig schon rathsam sein könnte, die sphäroidische Gestalt der Erde zu berücksichtigen, was zwar an sich nicht schwer sein, aber die Einfachheit der Rechnungen ohne allen Nutzen\par\clearpage
\SourceLeaf{143}\GaussVTheoryEarlyFolio{139}
\noindent sehr beeinträchtigen würde. Wir werden daher hier bei den zuletzt angeführten Formeln stehen bleiben, indem wir die Erde wie eine Kugel betrachten, deren Halbmesser $=R$ ist.

\Article{15}
Ist $X$ durch eine gegebene Function von $u$ und $\lambda$ ausgedrückt, so lässt sich daraus $Y$ a priori ableiten. Man setze das Integral $\int_0^u X\mathrm{d}u=T$, indem man bei der Integration $\lambda$ wie constant betrachtet; offenbar wird dann, wenn man auf gleiche Weise nach $u$ differentiirt, $\frac{\mathrm{d}(V+RT)}{\mathrm{d}u}=0$, mithin $V+RT$ eine von $u$ unabhängige Grösse, oder was dasselbe ist, in allen Punkten Eines Meridians constant; sie muss daher auch absolut constant sein, weil alle Meridiane in den Polen zusammenlaufen. Setzt man den Werth von $V$ im Nordpole $=V^*$, so wird also
\[
T=\frac{V^*-V}{R}
\]
und daher
\[
Y=\frac{\mathrm{d}T}{\sin u\cdot\mathrm{d}\lambda}
\]
Man kann dieses Resultat auch so ausdrücken:
\[
Y=\frac{1}{\sin u}\int_0^u\frac{\mathrm{d}X}{\mathrm{d}\lambda}\cdot\mathrm{d}u
\]

\Article{16}
Dieser merkwürdige Satz, dass \textit{wenn die nach Norden gerichtete Componente der horizontalen magnetischen Kraft für die ganze Erdoberfläche gegeben ist, die nach Westen (oder Osten) gerichtete Componente von selbst daraus folgt}, gilt verkehrt nur mit einer Modification. Ist nemlich $Y$ durch eine gegebene Function von $u$ und $\lambda$ ausgedrückt, und bezeichnet man mit $U$ das unbestimmte Integral $\int\sin u\cdot Y\mathrm{d}\lambda$, bei der Integration $u$ als constant angenommen, so wird $\frac{\mathrm{d}(V+RU)}{\mathrm{d}\lambda}=0$. oder $V+RU$ eine von $\lambda$ unabhängige Grösse, mithin allgemein zu reden eine Function von $u$. Es ist also auch $\frac{\mathrm{d}(V+RU)}{R\mathrm{d}u}=\frac{\mathrm{d}U}{\mathrm{d}u}-X$ eine solche Function, d. i. die Formel $\frac{\mathrm{d}U}{\mathrm{d}u}$ gibt einen unvollständigen Ausdruck von $X$, indem ein bloss $u$ enthaltender Bestandtheil unbestimmt bleibt. Dieser Mangel wird sich aber ergänzen lassen, wenn man ausser dem Ausdrucke für $Y$ auch den für $X$ in irgend Einem bestimmten Meridian, oder noch allgemeiner in irgend einer vom\par\GaussVTheoryGathering{18\textsuperscript{*}}
\par\clearpage
\SourceLeaf{144}\GaussVTheoryEarlyFolio{140}
\noindent Nordpol zum Südpol reichenden Linie besitzt. Man sieht also, dass, \textit{wenn man die Componente der horizontalen magnetischen Kraft in der Richtung nach Westen für die ganze Erdoberfläche, und die Componente in der Richtung nach Norden für alle Punkte in irgend einer vom Nordpol zum Südpol gehenden Linie kennt, die letztere Componente für die ganze Erdfläche von selbst daraus folgt.}

\Article{17}
Die vorhergehenden Untersuchungen beziehen sich allein auf den horizontalen Theil der erdmagnetischen Kraft: um auch den verticalen zu umfassen, müssen wir die Aufgabe in ihrer ganzen Allgemeinheit, also $V$ wie eine Function von \textit{dreien} veränderlichen Grössen betrachten, die den Platz eines unbestimmten Punktes im Raume $O$ ausdrücken. Wir wählen dazu die Entfernung $r$ vom Mittelpunkte der Erde, den Winkel $u$, welchen $r$ mit dem nordlichen Theile der Erdaxe macht, und den Winkel $\lambda$ zwischen der durch $r$ und die Erdaxe gelegten Ebene und einem festen Meridian, nach Osten zu als positiv gezählt.

Es sei die Function $V$ in eine nach den Potenzen von $r$ fallende Reihe entwickelt, der wir folgende Form geben
\[
V=\frac{RRP^0}{r}+\frac{R^3P'}{rr}+\frac{R^4P''}{r^3}+\frac{R^5P'''}{r^4}+\text{u. s. w.}
\]
Die Coëfficienten $P^0,P',P''$ u. s. w. sind hier Functionen von $u$ und $\lambda$; um zu übersehen, wie sie mit der Vertheilung des magnetischen Fluidums im Innern der Erde zusammenhängen, sei $\mathrm{d}\mu$ ein Element desselben, $\rho$ seine Entfernung von $O$, und für $\mathrm{d}\mu$ bedeuten $r^0,u^0,\lambda^0$ dasselbe, was $r,u,\lambda$ für $O$ sind. Man hat also $V=-\int\frac{\mathrm{d}\mu}{\rho}$ durch alle $\mathrm{d}\mu$ ausgedehnt; ferner
\[
\rho=\BareRoot{\{rr-2rr^0(\cos u\cos u^0+\sin u\sin u^0\cos(\lambda-\lambda^0))+r^0r^0\}}
\]
und wenn man $\frac1\rho$ in die Reihe entwickelt
\[
\frac1\rho=\frac1r(T^0+T'\cdot\frac{r^0}{r}+T''\cdot\frac{r^0r^0}{rr}+\text{u. s. w.})
\]
so wird
\[
\begin{aligned}RRP^0&=-\int T^0\mathrm{d}\mu,\quad R^3P'=-\int T'r^0\mathrm{d}\mu,\\R^4P''&=-\int T''r^0r^0\mathrm{d}\mu\quad\text{u. s. w.}\end{aligned}
\]
Da $T^0=1$ ist, so wird vermöge der Fundamentalvoraussetzung, dass die Menge des positiven und negativen Fluidums in jedem messbaren Theilchen seines Trägers, mithin auch in der \textit{ganzen} Erde, gleich gross, oder dass $\int\mathrm{d}\mu=0$ ist,
\par\clearpage
\SourceLeaf{145}\GaussVTheoryEarlyFolio{141}
\[
P^0=0
\]
oder das erste Glied unsrer Reihe für $V$ fällt aus.

Man sieht ferner, dass $P'$ die Form hat
\[
R^3P'=\alpha\cos u+\beta\sin u\cos\lambda+\gamma\sin u\sin\lambda
\]
wo $\alpha=-\int r^0\cos u^0\cdot\mathrm{d}\mu$, $\beta=-\int r^0\sin u^0\cos\lambda^0\mathrm{d}\mu$, $\gamma=-\int r^0\sin u^0\sin\lambda^0\cdot\mathrm{d}\mu$. Es sind also $-\alpha,-\beta,-\gamma$ nach der in der \textit{Intensitas vis magneticae} Art. 5 festgesetzten Erklärung die Momente des Erdmagnetismus in Beziehung auf drei rechtwinklige Axen, wovon die erste die Erdaxe, die zweite und dritte die Aequatorsradien für die Länge $0$ und $90^0$ sind.

Die allgemeinen Formeln für alle Coëfficienten der Reihe für $\frac1\rho$ können wir als bekannt voraussetzen; für unsern Zweck ist aber bloss nöthig zu bemerken, dass in Beziehung auf $u$ und $\lambda$ die Coëfficienten rationale ganze Functionen von $\cos u$, $\sin u\cos\lambda$ und $\sin u\sin\lambda$ sind, und zwar $T''$ von der zweiten Ordnung, $T'''$ von der dritten u. s. w. Dasselbe gilt also auch für die Coëfficienten $P'',P'''$ u. s. w.

Die Reihen für $\frac1\rho$ und für $V$ convergiren, solange $r$ nicht kleiner als $R$ ist, oder vielmehr, nicht kleiner, als der Halbmesser einer Kugel, welche die sämmtlichen magnetischen Theile der Erde einschliesst.

\Article{18}
Die Function $V$ thut, in Folge ihrer Zusammensetzung aus $-\int\frac{\mathrm{d}\mu}{\rho}$, folgender partiellen Differentialgleichung Genüge:
\[
0=\frac{r\mathrm{d}\mathrm{d}rV}{\mathrm{d}r^2}+\frac{\mathrm{d}\mathrm{d}V}{\mathrm{d}u^2}+\operatorname{cotg}u\cdot\frac{\mathrm{d}V}{\mathrm{d}u}+\frac1{\sin u^2}\cdot\frac{\mathrm{d}\mathrm{d}V}{\mathrm{d}\lambda^2}
\]
welche nichts anderes ist, als eine Umformung der bekannten
\[
0=\frac{\mathrm{d}\mathrm{d}V}{\mathrm{d}x^2}+\frac{\mathrm{d}\mathrm{d}V}{\mathrm{d}y^2}+\frac{\mathrm{d}\mathrm{d}V}{\mathrm{d}z^2}
\]
wo $x,y,z$ die rechtwinkligen Coordinaten von $O$ bedeuten. Substituirt man in jener den Werth von $V$
\[
V=\frac{R^3P'}{rr}+\frac{R^4P''}{r^3}+\frac{R^5P'''}{r^4}+\text{u. s. w.}
\]
so erhellt, dass für die einzelnen Coëfficienten $P',P'',P'''$ u. s. w. gleichfalls\par\clearpage
\SourceLeaf{146}\GaussVTheoryEarlyFolio{142}
\noindent partielle Differentialgleichungen Statt finden, deren allgemeiner Ausdruck ist
\[
0=n(n+1)P^{\mathrm{n}}+\frac{\mathrm{d}\mathrm{d}P^{\mathrm{n}}}{\mathrm{d}u^2}+\operatorname{cotg}u\frac{\mathrm{d}P^{\mathrm{n}}}{\mathrm{d}u}+\frac1{\sin u^2}\cdot\frac{\mathrm{d}\mathrm{d}P^{\mathrm{n}}}{\mathrm{d}\lambda^2}
\]

Aus dieser Gleichung, verbunden mit der Bemerkung im vorhergehenden Artikel, ergibt sich die allgemeine Form von $P^{\mathrm{n}}$. Bezeichnet man nemlich mit $P^{\mathrm{n,m}}$ folgende Function von $u$
\[
\begin{aligned}
\{\cos u^{n-m}&-\frac{(n-m)(n-m-1)}{2(2n-1)}\cos u^{n-m-2}\\
&+\frac{(n-m)(n-m-1)(n-m-2)(n-m-3)}{2\cdot4(2n-1)(2n-3)}\cos u^{n-m-4}\\&-\text{u. s. w.}\}\sin u^m
\end{aligned}
\]
so hat $P^{\mathrm{n}}$ die Form eines Aggregats von $2n+1$ Theilen
\[
\begin{aligned}
P^{\mathrm{n}}&=g^{\mathrm{n,0}}P^{\mathrm{n,0}}+(g^{\mathrm{n,1}}\cos\lambda+h^{\mathrm{n,1}}\sin\lambda)P^{\mathrm{n,1}}\\
&\quad+(g^{\mathrm{n,2}}\cos2\lambda+h^{\mathrm{n,2}}\sin2\lambda)P^{\mathrm{n,2}}+\text{etc.}+(g^{\mathrm{n,n}}\cos n\lambda+h^{\mathrm{n,n}}\sin n\lambda)P^{\mathrm{n,n}}
\end{aligned}
\]
wo $g^{\mathrm{n,0}},g^{\mathrm{n,1}},h^{\mathrm{n,1}},g^{\mathrm{n,2}}$ u. s. w. bestimmte Zahlcoëfficienten sind.

\Article{19}
Zerlegt man die in dem Punkte $O$ Statt findende magnetische Kraft in drei auf einander senkrechte $X,Y,Z$, wovon die dritte gegen den Mittelpunkt der Erde gerichtet ist, $X$ und $Y$ also die durch $O$ gelegte mit der Erde concentrische Kugelfläche berühren, und zwar $X$ in der durch $O$ und die Erdaxe gelegten Ebene nach Norden, $Y$ parallel mit dem Erdäquator nach Westen, so wird
\[
X=-\frac{\mathrm{d}V}{r\mathrm{d}u},\quad Y=-\frac{\mathrm{d}V}{r\sin u\,\mathrm{d}\lambda},\quad Z=-\frac{\mathrm{d}V}{\mathrm{d}r}
\]
folglich
\[
\begin{aligned}
X&=-\frac{R^3}{r^3}\left(\frac{\mathrm{d}P'}{\mathrm{d}u}+\frac Rr\cdot\frac{\mathrm{d}P''}{\mathrm{d}u}+\frac{RR}{rr}\cdot\frac{\mathrm{d}P'''}{\mathrm{d}u}+\text{u. s. w.}\right)\\
Y&=-\frac{R^3}{r^3\sin u}\left(\frac{\mathrm{d}P'}{\mathrm{d}\lambda}+\frac Rr\cdot\frac{\mathrm{d}P''}{\mathrm{d}\lambda}+\frac{RR}{rr}\cdot\frac{\mathrm{d}P'''}{\mathrm{d}\lambda}+\text{u. s. w.}\right)\\
Z&=\frac{R^3}{r^3}\left(2P'+\frac{3RP''}{r}+\frac{4RRP'''}{rr}+\text{u. s. w.}\right)
\end{aligned}
\]

Auf der Oberfläche der Erde sind $X,Y$ dieselben horizontalen Componenten, welche oben mit diesen Buchstaben bezeichnet sind, und $Z$ ist die verticale, positiv, wenn nach unten gerichtet. Die Ausdrücke für diese Kräfte auf der Oberfläche der Erde sind also
